\begin{scholium}\label{sga3:XXI:3.3.11}
Let \(P'\) be a set of roots satisfying (P 2) and closed. There exist a
system of simple roots \(\Delta\) and a subset \(\Delta'\) of
\(\Delta\) such that
\[
 P'=R\cap\bigl(\mathbb{N}\cdot\Delta
          \,\cup\, -\mathbb{N}\cdot\Delta'\bigr).
\]
If
\[
 R_{\Delta'}=(\mathbb{Z}\cdot\Delta')\cap R,
\]
which is a closed and symmetric subset of \(R\), then \(P'\) is the
disjoint union of \(R_{\Delta'}\) and the closed subset of
\(\mathscr{P}(\Delta)\) formed by the
\[
 \alpha\in\mathscr{P}(\Delta)\setminus
           \mathbb{N}\cdot\Delta',
\]
that is, the positive roots whose decomposition with respect to
\(\Delta\) ``contains at least one root of
\(\Delta\setminus\Delta'\)''.
\end{scholium}

\subsection{Closed and symmetric sets of roots}
\label{sga3:XXI:3.4}

\begin{proposition}\label{sga3:XXI:3.4.1}
Let
\[
 \cR=(M,M^\ast,R,R^\ast)
\]
be a root datum, and let \(R'\) be a closed and symmetric subset of
\(R\). Then:
\begin{enumerate}
\item[(i)] \(\cR'=(M,M^\ast,R',R'^\ast)\) is a root datum;
\item[(ii)] for every system of positive roots \(R_+\) of \(R\),
  \(R'_+=R_+\cap R'\) is a system of positive roots of \(R'\);
\item[(iii)] the Weyl group \(W(\cR')\) of \(\cR'\) is the subgroup
  of \(W(\cR)\) generated by the \(s_\alpha\), for
  \(\alpha\in R'\).
\end{enumerate}
\end{proposition}

\begin{proof}
The first assertion is immediate from
Corollary~\ref{sga3:XXI:1.2.11}; the second follows from
Corollary~\ref{sga3:XXI:3.3.8}; and the third is evident.
\end{proof}

\begin{corollary}\label{sga3:XXI:3.4.2}
Let \(w\in W(\cR')\). The order of \(w\) is the smallest integer
\(n>0\) such that
\[
 w^n(\alpha')=\alpha'
 \qquad\text{for every }\alpha'\in R'.
\]
\end{corollary}

\begin{proof}
Apply Corollary~\ref{sga3:XXI:1.2.7} to the root datum \(\cR'\).
\end{proof}

\begin{corollary}\label{sga3:XXI:3.4.3}
Let \(\alpha\) and \(\beta\) be two nonproportional roots. Let \(n\)
be the smallest integer \(>0\) such that
\[
 (s_\alpha s_\beta)^n(\alpha)=\alpha,
 \qquad
 (s_\alpha s_\beta)^n(\beta)=\beta.
\]
Then the subgroup \(W_{\alpha,\beta}\) of \(W\) generated by
\(s_\alpha\) and \(s_\beta\) is defined by the relations
\[
 s_\alpha^2=1,\qquad
 s_\beta^2=1,\qquad
 (s_\alpha s_\beta)^n=1.
\]
\end{corollary}

\begin{proof}
In view of Corollary~\ref{sga3:XXI:3.4.2}, it remains only to verify
the following lemma.
\end{proof}

\begin{lemma}\label{sga3:XXI:3.4.4}
Let \(G\) be the group generated by two elements \(x\) and \(y\)
subject to the relations \(x^2=y^2=1\). Every normal subgroup of
\(G\) containing neither \(x\) nor \(y\) is generated, as a normal
subgroup, by an element of the form \((xy)^n\), where \(n>0\).
\end{lemma}

\begin{proof}
\footnote{Editors' note: The original has been corrected in what
follows.}
Every element of \(G\) can be written as \((xy)^n\), as
\((yx)^n=(xy)^{-n}\), or in one of the forms
\[
\begin{array}{c@{\qquad}c@{\qquad}c}
 (a) & x(yx)^{2n} & \text{or } y(xy)^{2n+1},\\[2mm]
 (b) & y(xy)^{2n} & \text{or } x(yx)^{2n+1},
\end{array}
\]
where \(n\in\mathbb{N}\). The elements of type (a), respectively of
type (b), are conjugate to \(x\), respectively to \(y\).
\end{proof}

\begin{remark}\label{sga3:XXI:3.4.5}
The integer \(n\) is computed immediately. Put
\[
 (\alpha^\ast,\beta)=p,\qquad
 (\beta^\ast,\alpha)=q.
\]
Then\footnote{Editors' note: The original has been corrected in what
follows.}
\[
 (s_\alpha s_\beta)(\alpha)=(pq-1)\alpha-q\beta,
 \qquad
 (s_\alpha s_\beta)(\beta)=p\alpha-\beta.
\]
Thus \(n\) is the order of the matrix
\[
 \begin{pmatrix}
  pq-1 & p\\
  -q   & -1
 \end{pmatrix}.
\]
If \(pq=0\), then \(p=q=0\) by
Corollary~\ref{sga3:XXI:2.2.2}, whence \(n=2\). Otherwise,
Proposition~\ref{sga3:XXI:2.3.1} shows that \(pq\) equals \(1\),
\(2\), or \(3\), and one finds respectively \(n=3\), \(4\), or
\(6\).

\smallskip
\noindent
\emph{N.B.} Expressing the finiteness of the order of the preceding
matrix recovers inequality \SGARefLink{sga3:XXI:eq:13}{(13)} of
Proposition~\ref{sga3:XXI:2.3.1}.
\end{remark}
